\section{模型评价}\label{sec:model_evaluation}

\subsection{模型的主要特点}\label{subsec:evaluation_features}

本文方法在统一时间位置之外保留了输入来源。问题一的每个有效公共窗记录模态掩码、原生索引、时间支持以及文本、WAV或源帧信息。公共位置在形成定长表示的同时仍保留原始来源；后续检查异常输入或解释具体位置时，可以沿来源记录返回原始观测。无有效交叠的位置则用掩码与空来源明确区分。

预测部分使用规模较小的双任务结构。MMP给出等权聚合参照，LTARP只增加三个模态的内容打分器和残差系数，在原有163460个参数之外增加883个可训练参数，总量为164343。新增路径依当前输入调整时间位置权重，与均值表示共同进入既有融合和输出层。这一结构便于把性能差异与池化方式联系起来，但参数增加较少本身并不保证每项指标都改善。

连续缺失评价同时改变模态、比例和位置，并在相同受损输入上计算极性分类与强度回归指标。完整输入与54种固定场景分别记录，再按既定口径汇总，因此可以查看平均表现，也可以辨别哪些输入受损时结果变化较大。解释部分则把模态贡献、联合输入作用和局部遮挡放在固定模型下计算；关键索引能够进一步与已有来源记录对应，而不只停留在总体重要性排序。

\subsection{稳定性与适用性分析}\label{subsec:evaluation_stability}

现有实验从不同层面检查了结果。问题一对对齐与视觉编码方案使用5次重复分层五折探针，作用是比较同一批100条视频上的特征方案。问题二使用三个随机初始化训练模型；对每个随机种子固定已训练参数后，在完整输入和54种连续缺失场景上统一评估，因此能够比较模态、比例和位置变化对同一模型配置的影响。问题三将附件2验证集按video ID分成窗口比例设计子集与独立验证子集，以随机等长窗口为局部位置对照，并按视频分组自助采样给出区间。这些设计分别针对方案选择、预测波动和解释位置的模型响应，不能互相替代。

稳定性并不在所有层面相同。LTARP相对MMP的宏平均F1、MAE和Pearson总体改善较小，准确率在部分比较中下降，三次初始化的综合增益也有一次近乎持平；因此不宜把局部收益表述为所有指标的共同提升。缺失文本的前、中、后段对分类和回归产生的变化并不一致，提示应用时须结合输出目标判断受损位置的影响。跨种子解释中，分类主导模态的一致率高于关键区间交并比所反映的边界稳定性，局部窗口应结合相邻位置的响应曲线阅读。上述适用性限于本文给定数据、缺失规则和固定预测模型。

\subsection{模型局限与改进方向}\label{subsec:evaluation_limits}

附件1仅有100条视频，特征方案的比较样本较少。问题一探针按片段标签进行重复分层五折，未按原始视频ID分组，结果适合内部方案筛选，不能直接视作跨视频泛化估计。后续若扩展同类数据，可增加视频层面的独立评估，并维持时间来源记录以检查不同采集条件。

问题二以规则化的连续置零模拟缺失，覆盖指定模态、位置和比例，却不能穷尽非规则的传感中断或模态质量退化。后续可在既有固定评价协议外加入不规则区间与非零噪声扰动，并分别报告结果，避免把不同故障机制混为一类。标签分布也不均衡，中性类别召回低于两侧极性类别；后续可针对中性与弱极性样本检查标注边界和分类校准，而不能仅凭总体准确率判断这一类别的改进。

问题三的贡献和遮挡效应都来自固定模型的输入扰动，表示模型输出对输入的响应，不等于现实情感形成的因果作用。局部区间的边界还会随随机初始化变化，适合与模态层面的结果及邻近窗口一并解释；后续可比较多种窗口宽度或将相邻高响应区间聚合。文本关键位置可回到原文字符，而附件4的语音与视觉目前主要回到未对齐特征行，来源精度尚不相同。若补足这些附件的媒体映射记录，可将后两种模态进一步连接到真实音频区间和视频帧。

总体上，本文完成了原始三模态特征构建、连续缺失条件下的情感预测，以及模态与关键位置的模型解释，并给出附件3和附件4的无标签样本输出。
